%%%%%%%%%%%%%%%%%%%%%%%%%%%%%%%%%%%%%%%%%%%%%%%%%%%%%%%%%%%%%%%%%%%%%%%%%%%%%%%
%%  ABSCHLUSSBERICHT: Pflegegrad-Analyse & Höherstufungsstrategie
%%  Session: EBF-S-2026-01-25-HLT-001
%%  Datum: 25. Januar 2026
%%
%%  Kompilieren mit: pdflatex Pflegegrad_Analyse_Bericht.tex
%%%%%%%%%%%%%%%%%%%%%%%%%%%%%%%%%%%%%%%%%%%%%%%%%%%%%%%%%%%%%%%%%%%%%%%%%%%%%%%

\documentclass[11pt,a4paper]{article}

% Pakete
\usepackage[utf8]{inputenc}
\usepackage[T1]{fontenc}
\usepackage[ngerman]{babel}
\usepackage{geometry}
\usepackage{graphicx}
\usepackage{xcolor}
\usepackage{booktabs}
\usepackage{longtable}
\usepackage{array}
\usepackage{multirow}
\usepackage{fancyhdr}
\usepackage{titlesec}
\usepackage{tcolorbox}
\usepackage{enumitem}
\usepackage{hyperref}
\usepackage{eurosym}
\usepackage{amssymb}
\usepackage{pifont}

% Seitenränder
\geometry{left=2.5cm, right=2.5cm, top=2.5cm, bottom=2.5cm}

% Farben
\definecolor{primaryblue}{RGB}{0, 82, 147}
\definecolor{successgreen}{RGB}{40, 167, 69}
\definecolor{warningorange}{RGB}{255, 152, 0}
\definecolor{dangerred}{RGB}{220, 53, 69}
\definecolor{lightgray}{RGB}{248, 249, 250}
\definecolor{boxblue}{RGB}{232, 244, 253}
\definecolor{boxgreen}{RGB}{232, 245, 233}
\definecolor{boxorange}{RGB}{255, 243, 224}
\definecolor{boxred}{RGB}{255, 235, 238}

% Checkmark und X
\newcommand{\cmark}{\ding{51}}
\newcommand{\xmark}{\ding{55}}

% Boxen
\tcbuselibrary{skins,breakable}

\newtcolorbox{infobox}[1][]{
    colback=boxblue,
    colframe=primaryblue,
    fonttitle=\bfseries,
    title=#1,
    breakable
}

\newtcolorbox{successbox}[1][]{
    colback=boxgreen,
    colframe=successgreen,
    fonttitle=\bfseries,
    title=#1,
    breakable
}

\newtcolorbox{warningbox}[1][]{
    colback=boxorange,
    colframe=warningorange,
    fonttitle=\bfseries,
    title=#1,
    breakable
}

\newtcolorbox{dangerbox}[1][]{
    colback=boxred,
    colframe=dangerred,
    fonttitle=\bfseries,
    title=#1,
    breakable
}

\newtcolorbox{graybox}{
    colback=lightgray,
    colframe=gray,
    breakable
}

% Kopf- und Fußzeile
\pagestyle{fancy}
\fancyhf{}
\fancyhead[L]{\small Pflegegrad-Analyse}
\fancyhead[R]{\small EBF-S-2026-01-25-HLT-001}
\fancyfoot[C]{\thepage}
\renewcommand{\headrulewidth}{0.4pt}

% Titel-Formatierung
\titleformat{\section}{\Large\bfseries\color{primaryblue}}{\thesection}{1em}{}
\titleformat{\subsection}{\large\bfseries\color{primaryblue}}{\thesubsection}{1em}{}

% Hyperref Setup
\hypersetup{
    colorlinks=true,
    linkcolor=primaryblue,
    urlcolor=primaryblue,
    pdftitle={Pflegegrad-Analyse und Höherstufungsstrategie},
    pdfauthor={EBF Framework}
}

%%%%%%%%%%%%%%%%%%%%%%%%%%%%%%%%%%%%%%%%%%%%%%%%%%%%%%%%%%%%%%%%%%%%%%%%%%%%%%%
\begin{document}

%% Titelseite
\begin{titlepage}
    \centering
    \vspace*{2cm}

    {\Huge\bfseries\color{primaryblue} Abschlussbericht}\\[0.5cm]
    {\LARGE Pflegegrad-Analyse \& Höherstufungsstrategie}\\[2cm]

    \begin{tcolorbox}[colback=lightgray, colframe=gray, width=12cm]
        \centering
        \begin{tabular}{ll}
            \textbf{Session:} & EBF-S-2026-01-25-HLT-001 \\
            \textbf{Datum:} & 25. Januar 2026 \\
            \textbf{Modus:} & STANDARD \\
            \textbf{Domain:} & Health (HLT) \\
        \end{tabular}
    \end{tcolorbox}

    \vspace{2cm}

    \begin{infobox}[Kernempfehlung]
        \begin{enumerate}
            \item \textbf{SOFORT} Höherstufungsantrag stellen $\rightarrow$ Ziel: Pflegegrad 4
            \item \textbf{Parallel} auf Pflegeheim-Warteliste setzen (12--18 Monate!)
            \item \textbf{24h-Pflege} als Brückenlösung prüfen ($\sim$1.400\euro{}/Monat netto)
            \item \textbf{Verhinderungspflege} nutzen $\rightarrow$ Ehemann entlasten
        \end{enumerate}
    \end{infobox}

    \vspace{2cm}

    {\large\textbf{Finanzieller Mehrwert bei Höherstufung PG2 $\rightarrow$ PG4:}}\\[0.5cm]
    {\Large\color{successgreen}\textbf{5.200 -- 12.200 \euro{}/Jahr mehr Leistungen!}}

    \vfill

    {\small Erstellt mit dem Evidence-Based Framework (EBF)}

\end{titlepage}

%% Inhaltsverzeichnis
\tableofcontents
\newpage

%%%%%%%%%%%%%%%%%%%%%%%%%%%%%%%%%%%%%%%%%%%%%%%%%%%%%%%%%%%%%%%%%%%%%%%%%%%%%%%
\section{Executive Summary}

\begin{successbox}[Zusammenfassung auf einen Blick]

\textbf{Ausgangslage:}
\begin{itemize}[noitemsep]
    \item 83-jährige Patientin mit mittelgradiger Demenz, Alzheimer-Anzeichen
    \item Aktuell Pflegegrad 2 -- \textbf{DEUTLICH ZU NIEDRIG}
    \item Ehemann (87J) pflegt alleine zuhause, ist überfordert
    \item 3 Töchter wollen Eltern nicht trennen
\end{itemize}

\vspace{0.5cm}
\textbf{Finanzieller Mehrwert bei Höherstufung PG2 $\rightarrow$ PG4:}
\begin{itemize}[noitemsep]
    \item +433\euro{}/Monat Pflegegeld (332 $\rightarrow$ 765\euro{})
    \item +1.017\euro{}/Monat Sachleistungen (761 $\rightarrow$ 1.778\euro{})
    \item = \textbf{5.200 -- 12.200\euro{}/Jahr mehr Leistungen!}
\end{itemize}

\vspace{0.5cm}
\textbf{Erfolgswahrscheinlichkeit:}
\begin{itemize}[noitemsep]
    \item PG3 oder höher: \textbf{85\%} (mit guter Vorbereitung)
    \item PG4: \textbf{65\%} (mit Pflegetagebuch + allen Interventionen)
\end{itemize}

\end{successbox}

%%%%%%%%%%%%%%%%%%%%%%%%%%%%%%%%%%%%%%%%%%%%%%%%%%%%%%%%%%%%%%%%%%%%%%%%%%%%%%%
\section{Fallbeschreibung}

\subsection{Patientin}

\begin{center}
\begin{tabular}{ll}
\toprule
\textbf{Merkmal} & \textbf{Beschreibung} \\
\midrule
Alter & 83 Jahre \\
Geschlecht & Weiblich \\
Wohnort & Hessen, Deutschland \\
Aktueller Pflegegrad & 2 \\
Hauptdiagnose & Mittelgradige Demenz mit Alzheimer-Anzeichen \\
Mobilität & Stark eingeschränkt, kann nur schlecht laufen \\
Sturzrisiko & Hoch, regelmäßige Stürze \\
Kognition & Kurzzeitgedächtnis mangelhaft \\
Vorerkrankung & Brustkrebs (über Jahre) \\
\bottomrule
\end{tabular}
\end{center}

\subsection{Familiäre Situation}

\begin{center}
\begin{tabular}{lll}
\toprule
\textbf{Person} & \textbf{Rolle} & \textbf{Situation} \\
\midrule
Ehemann & Hauptpflegeperson & 87 Jahre, will unbedingt zuhause bleiben \\
3 Töchter & Unterstützung & Wollen Eltern nicht trennen \\
\bottomrule
\end{tabular}
\end{center}

\subsection{Präferenzen der Familie}

\begin{itemize}
    \item[\cmark] Häusliche Pflege bevorzugt (Wunsch des Vaters)
    \item[\cmark] Eheleute sollen zusammenbleiben
    \item[$\triangle$] Pflegeheim als Absicherung in Diskussion
    \item[$\triangle$] Lange Wartezeiten für Demenz-Pflegeplätze bekannt
\end{itemize}

%%%%%%%%%%%%%%%%%%%%%%%%%%%%%%%%%%%%%%%%%%%%%%%%%%%%%%%%%%%%%%%%%%%%%%%%%%%%%%%
\section{Analyse: Warum Pflegegrad 2 zu niedrig ist}

\subsection{NBA-Modul-Analyse}

Das Neue Begutachtungsassessment (NBA) bewertet die Selbständigkeit in 6 Modulen:

\begin{center}
\begin{tabular}{lccl}
\toprule
\textbf{Modul} & \textbf{Geschätzt} & \textbf{Max} & \textbf{Begründung} \\
\midrule
1. Mobilität & 6--7,5 & 10 & Sturzgefahr, kann nur schlecht laufen \\
2/3. Kognition & 11--15 & 15 & Mittelgradige Demenz, Alzheimer \\
4. Selbstversorgung & 25--35 & 40 & Braucht Anleitung bei allem \\
5. Therapie/Krankheit & 5--15 & 20 & Medikamente, Brustkrebs-Nachsorge \\
6. Alltagsleben & 10--15 & 15 & Keine Eigenstruktur mehr möglich \\
\midrule
\textbf{GESAMT} & \textbf{57--87,5} & \textbf{100} & \\
\bottomrule
\end{tabular}
\end{center}

\begin{warningbox}[Pflegegrad-Schwellen]
\begin{itemize}[noitemsep]
    \item PG 2: 27,0 -- 47,4 Punkte $\leftarrow$ \textbf{AKTUELL (zu niedrig!)}
    \item PG 3: 47,5 -- 69,9 Punkte $\leftarrow$ \textbf{MINDESTENS realistisch}
    \item PG 4: 70,0 -- 89,9 Punkte $\leftarrow$ \textbf{WAHRSCHEINLICH bei guter Dokumentation}
\end{itemize}
\end{warningbox}

\subsection{Vermutete Gründe für Unterbewertung}

\begin{center}
\begin{tabular}{lcc}
\toprule
\textbf{Ursache} & \textbf{Häufigkeit} & \textbf{Wahrscheinlich?} \\
\midrule
``Guter Tag'' bei Begutachtung & Sehr häufig & \cmark~Ja \\
Angehörige spielen herunter & Sehr häufig & \cmark~Ja \\
Ehemann sagt ``wir schaffen das'' & Häufig & \cmark~Ja \\
Demenz wurde unterschätzt & Häufig & \cmark~Wahrscheinlich \\
Sturzrisiko nicht dokumentiert & Häufig & \cmark~Kein Protokoll \\
\bottomrule
\end{tabular}
\end{center}

%%%%%%%%%%%%%%%%%%%%%%%%%%%%%%%%%%%%%%%%%%%%%%%%%%%%%%%%%%%%%%%%%%%%%%%%%%%%%%%
\section{Finanzielle Analyse}

\subsection{Pflegekassen-Leistungen nach Pflegegrad}

\begin{center}
\begin{tabular}{lrrrc}
\toprule
\textbf{Leistung (monatlich)} & \textbf{PG 2} & \textbf{PG 3} & \textbf{PG 4} & \textbf{Differenz} \\
\midrule
Pflegegeld & 332\euro{} & 573\euro{} & 765\euro{} & +433\euro{} \\
Sachleistungen ambulant & 761\euro{} & 1.432\euro{} & 1.778\euro{} & +1.017\euro{} \\
Tagespflege & 689\euro{} & 1.298\euro{} & 1.612\euro{} & +923\euro{} \\
Entlastungsbetrag & 125\euro{} & 125\euro{} & 125\euro{} & --- \\
\midrule
Entlastungsbudget (jährlich) & 3.539\euro{} & 3.539\euro{} & 3.539\euro{} & --- \\
Wohnraumanpassung (einmalig) & 4.000\euro{} & 4.000\euro{} & 4.000\euro{} & --- \\
\bottomrule
\end{tabular}
\end{center}

\subsection{Pflegeheim Hessen: Eigenanteil (Januar 2026)}

\begin{infobox}[Eigenanteil Pflegeheim Hessen]
\textbf{Gesamtkosten (brutto):} $\sim$5.300\euro{}/Monat

\vspace{0.3cm}
\textbf{Eigenanteil nach Aufenthaltsdauer:}
\begin{itemize}[noitemsep]
    \item 1. Jahr: 3.229\euro{}/Monat
    \item 2. Jahr: 2.924\euro{}/Monat
    \item 3. Jahr: $\sim$2.500\euro{}/Monat
    \item Ab 4. Jahr: 2.011\euro{}/Monat
\end{itemize}

\vspace{0.3cm}
\textbf{Aufschlüsselung:}
\begin{itemize}[noitemsep]
    \item Pflegekosten: 2.030\euro{} (inkl. 161\euro{} Ausbildung)
    \item Unterkunft \& Verpflegung: 978\euro{}
    \item Investitionskosten: 525\euro{}
\end{itemize}

\vspace{0.3cm}
\textbf{Wartezeit Demenz-Spezialbereich: 12--18 Monate}
\end{infobox}

\subsection{Alternative: 24-Stunden-Pflege zuhause}

\begin{successbox}[24h-Pflege: Kostenrechnung bei Pflegegrad 4]

\textbf{Kosten (osteuropäische Betreuungskraft):}
\begin{itemize}[noitemsep]
    \item Basisbetreuung: 2.200 -- 2.500\euro{}/Monat
    \item Bei PG4/höherem Bedarf: 2.550 -- 2.890\euro{}/Monat
\end{itemize}

\vspace{0.3cm}
\textbf{Gegenfinanzierung:}
\begin{itemize}[noitemsep]
    \item Pflegegeld PG4: $-$765\euro{}
    \item Steuererstattung (20\%): $-$333\euro{}/Monat
    \item Verhinderungspflege anteilig: $-$295\euro{}/Monat
\end{itemize}

\vspace{0.3cm}
\textbf{NETTO-EIGENANTEIL: ca. 1.300 -- 1.500\euro{}/Monat}

\vspace{0.3cm}
\textbf{Vergleich:}
\begin{itemize}[noitemsep]
    \item Pflegeheim (1. Jahr): $\sim$3.229\euro{}/Monat
    \item 24h-Pflege (netto): $\sim$1.400\euro{}/Monat
    \item \textbf{ERSPARNIS: $\sim$1.800\euro{}/Monat = 21.600\euro{}/Jahr}
\end{itemize}

\vspace{0.3cm}
\textbf{Vorteile:}
\begin{itemize}[noitemsep]
    \item[\cmark] Beide Eheleute bleiben zusammen
    \item[\cmark] Sofort verfügbar (5--7 Tage)
    \item[\cmark] Günstiger als Pflegeheim
    \item[\cmark] Ehemann (87J) wird entlastet
\end{itemize}

\end{successbox}

\subsection{Kombinationsleistung (Beispielrechnung PG4)}

\begin{graybox}
\textbf{Budgets bei PG4:}
\begin{itemize}[noitemsep]
    \item Pflegegeld 100\%: 800\euro{}/Monat
    \item Sachleistungen 100\%: 1.859\euro{}/Monat
\end{itemize}

\vspace{0.3cm}
\textbf{Beispiel: 40\% Sachleistung + 60\% Pflegegeld}
\begin{itemize}[noitemsep]
    \item Sachleistung: 40\% $\times$ 1.859\euro{} = 743,60\euro{}
    \item Pflegegeld: 60\% $\times$ 800\euro{} = 480,00\euro{}
    \item \textbf{GESAMT: 1.223,60\euro{}}
\end{itemize}

\vspace{0.3cm}
\textbf{PLUS (kürzt Kombileistung NICHT):}
\begin{itemize}[noitemsep]
    \item Tagespflege 3$\times$/Woche: $\sim$800\euro{}
    \item Entlastungsbetrag: 125\euro{}
    \item Entlastungsbudget anteilig: $\sim$295\euro{}
\end{itemize}

\textbf{Monatliche Leistungen gesamt: $\sim$2.450\euro{}}
\end{graybox}

%%%%%%%%%%%%%%%%%%%%%%%%%%%%%%%%%%%%%%%%%%%%%%%%%%%%%%%%%%%%%%%%%%%%%%%%%%%%%%%
\section{Handlungsplan: 4 Phasen}

\subsection{Phase 1: SOFORT (Diese Woche)}

\begin{center}
\begin{tabular}{clll}
\toprule
\textbf{Prio} & \textbf{Aufgabe} & \textbf{Verantwortlich} & \textbf{Frist} \\
\midrule
\textcolor{dangerred}{\textbf{!}} & Höherstufungsantrag stellen & Schwester 1 & Montag \\
\textcolor{dangerred}{\textbf{!}} & Sturzprotokoll beginnen & Alle & Sofort \\
\textcolor{warningorange}{\textbf{!}} & Pflegeheim-Warteliste (3--5 Heime) & Schwester 3 & Diese Woche \\
\textcolor{warningorange}{\textbf{!}} & Hausnotruf installieren & Schwester 2 & Diese Woche \\
\bottomrule
\end{tabular}
\end{center}

\subsection{Phase 2: VORBEREITUNG (2--4 Wochen)}

\begin{center}
\begin{tabular}{clll}
\toprule
\textbf{Prio} & \textbf{Aufgabe} & \textbf{Verantwortlich} & \textbf{Frist} \\
\midrule
\textcolor{dangerred}{\textbf{!}} & Pflegetagebuch täglich führen & Alle/Ehemann & Täglich \\
\textcolor{dangerred}{\textbf{!}} & Ärztliche Unterlagen sammeln & Schwester 2 & 2 Wochen \\
\textcolor{warningorange}{\textbf{!}} & 24h-Pflege-Angebote einholen & Schwester 3 & 2 Wochen \\
\textcolor{warningorange}{\textbf{!}} & Schwestern briefen für MDK & Schwester 1 & Vor Termin \\
\textcolor{warningorange}{\textbf{!}} & Ehemann vorbereiten (Reframing) & Alle & Vor Termin \\
\bottomrule
\end{tabular}
\end{center}

\subsection{Phase 3: MDK-BEGUTACHTUNG (Woche 3--6)}

\begin{center}
\begin{tabular}{cll}
\toprule
\textbf{Prio} & \textbf{Aufgabe} & \textbf{Verantwortlich} \\
\midrule
\textcolor{dangerred}{\textbf{!}} & Mindestens 1 Schwester dabei & Schwester 1 oder 2 \\
\textcolor{dangerred}{\textbf{!}} & Pflegetagebuch vorzeigen & Wer dabei ist \\
\textcolor{dangerred}{\textbf{!}} & Mutter korrigieren wenn nötig & Wer dabei ist \\
\textcolor{dangerred}{\textbf{!}} & Eigenes Protokoll anfertigen & Wer dabei ist \\
\bottomrule
\end{tabular}
\end{center}

\subsection{Phase 4: NACH BESCHEID (Woche 6--10+)}

\begin{center}
\begin{tabular}{lll}
\toprule
\textbf{Szenario} & \textbf{Aktion} & \textbf{Frist} \\
\midrule
\textcolor{successgreen}{\cmark} PG3/4 bewilligt & Leistungen einrichten, 24h-Pflege starten & Sofort \\
\textcolor{warningorange}{$\triangle$} Nur PG3 & Gutachten prüfen, ggf. Widerspruch & 1 Monat \\
\textcolor{dangerred}{\xmark} Bleibt PG2 & \textbf{WIDERSPRUCH einlegen!} & 1 Monat \\
\bottomrule
\end{tabular}
\end{center}

%%%%%%%%%%%%%%%%%%%%%%%%%%%%%%%%%%%%%%%%%%%%%%%%%%%%%%%%%%%%%%%%%%%%%%%%%%%%%%%
\section{Verhaltensinterventionen für die Familie}

\subsection{Übersicht}

\begin{center}
\begin{tabular}{clll}
\toprule
\textbf{\#} & \textbf{Intervention} & \textbf{Ziel} & \textbf{Wirkung} \\
\midrule
1 & Kognitives Reframing & Scham/Stolz überwinden & \ding{72}\ding{72}\ding{72}\ding{72}\ding{72} \\
2 & Rollen-Commitment & Last verteilen & \ding{72}\ding{72}\ding{72}\ding{72} \\
3 & Pre-Mortem MDK & Fehler vermeiden & \ding{72}\ding{72}\ding{72}\ding{72}\ding{72} \\
4 & Defaults für Ehemann & Widerstand lösen & \ding{72}\ding{72}\ding{72}\ding{72} \\
5 & Wenn-Dann-Pläne & Umsetzung sichern & \ding{72}\ding{72}\ding{72}\ding{72} \\
\bottomrule
\end{tabular}
\end{center}

\subsection{Kernbotschaften}

\begin{infobox}[Für die Schwestern]
``Der MDK-Gutachter sieht 60 Minuten. Wir leben 24 Stunden täglich mit der Situation. Es ist unsere \textbf{PFLICHT}, die Realität zu zeigen -- für Mama UND für Papa.''
\end{infobox}

\begin{infobox}[Für den Ehemann (87J)]
``Hilfe anzunehmen ist \textbf{STARK}, nicht schwach. Du sorgst für Mama, indem du auf DICH achtest. Was wäre, wenn DIR etwas passiert?''
\end{infobox}

\begin{infobox}[Bei der MDK-Begutachtung]
``Wir haben 50 Jahre eingezahlt. Das ist unser \textbf{RECHT}, kein Almosen.''
\end{infobox}

\subsection{Rollen-Verteilung unter den Schwestern}

\begin{graybox}
\textbf{Schwester 1: ``Koordinatorin''}
\begin{itemize}[noitemsep]
    \item Hauptansprechpartnerin für Pflegekasse
    \item Schreibt Höherstufungsantrag
    \item Organisiert MDK-Termin
    \item Führt Pflegetagebuch (oder delegiert)
\end{itemize}

\textbf{Schwester 2: ``Dokumentarin''}
\begin{itemize}[noitemsep]
    \item Sammelt alle ärztlichen Unterlagen
    \item Führt Sturzprotokoll
    \item Fotografiert Hilfsmittel/Wohnsituation
    \item Ist beim MDK-Termin dabei als Zeugin
\end{itemize}

\textbf{Schwester 3: ``Entlasterin''}
\begin{itemize}[noitemsep]
    \item Kümmert sich um Papa (87J) -- ER braucht auch Fürsorge!
    \item Recherchiert 24h-Pflege-Optionen
    \item Besucht Pflegeheime (Warteliste)
    \item Organisiert Verhinderungspflege
\end{itemize}
\end{graybox}

%%%%%%%%%%%%%%%%%%%%%%%%%%%%%%%%%%%%%%%%%%%%%%%%%%%%%%%%%%%%%%%%%%%%%%%%%%%%%%%
\section{Risiken und Absicherung}

\subsection{Risiko-Matrix}

\begin{center}
\begin{tabular}{lccc}
\toprule
\textbf{Risiko} & \textbf{Wahrsch.} & \textbf{Auswirkung} & \textbf{Maßnahme} \\
\midrule
Höherstufung abgelehnt & 15\% & Hoch & Widerspruch! \\
Ehemann (87J) fällt aus & Mittel & Kritisch & Notfallplan \\
Pflegeheim-Platz nicht frei & Hoch & Hoch & JETZT Warteliste \\
24h-Kraft nicht verfügbar & Gering & Mittel & Mehrere Agenturen \\
Mutter stürzt schwer & Hoch & Kritisch & Hausnotruf SOFORT \\
Familie uneinig & Gering & Mittel & Wöchentl. Call \\
\bottomrule
\end{tabular}
\end{center}

\subsection{Notfallplan: Wenn der Ehemann ausfällt}

\begin{dangerbox}[Notfallplan]
\textbf{WENN} der Ehemann (87J) ins Krankenhaus muss oder ausfällt:

\vspace{0.3cm}
\textbf{SOFORT (Tag 1):}
\begin{itemize}[noitemsep]
    \item[$\square$] Schwester springt ein (wer kann am schnellsten?)
    \item[$\square$] Pflegekasse anrufen: Kurzzeitpflege beantragen
    \item[$\square$] 24h-Pflege-Agentur anrufen: Notfall-Vermittlung (5--7 Tage)
\end{itemize}

\textbf{KURZFRISTIG (Woche 1):}
\begin{itemize}[noitemsep]
    \item[$\square$] Kurzzeitpflege-Platz suchen
    \item[$\square$] Pflegeheime kontaktieren: Notfall-Aufnahme möglich?
\end{itemize}

\textbf{VORBEREITUNG JETZT:}
\begin{itemize}[noitemsep]
    \item[$\square$] Telefonnummern aller Agenturen/Heime griffbereit
    \item[$\square$] Vollmachten für alle 3 Schwestern
    \item[$\square$] Medikamentenplan aktuell und lesbar
\end{itemize}
\end{dangerbox}

%%%%%%%%%%%%%%%%%%%%%%%%%%%%%%%%%%%%%%%%%%%%%%%%%%%%%%%%%%%%%%%%%%%%%%%%%%%%%%%
\section{Mustervorlagen}

\subsection{Höherstufungsantrag}

\begin{graybox}
\small
An: [Pflegekasse]\\
Datum: [heute]\\
Betreff: Antrag auf Höherstufung des Pflegegrades\\
Versicherte: [Name], geb. [Datum], Vers.-Nr.: [Nummer]

\vspace{0.3cm}
Sehr geehrte Damen und Herren,

hiermit beantrage ich für die oben genannte Versicherte die Höherstufung des Pflegegrades.

\textbf{Begründung:}\\
Der Gesundheitszustand hat sich seit der letzten Begutachtung wesentlich verschlechtert:

\begin{itemize}[noitemsep]
    \item Fortschreitende mittelgradige Demenz mit Alzheimer-Symptomen
    \item Stark eingeschränktes Kurzzeitgedächtnis
    \item Erhebliche Mobilitätseinschränkungen
    \item Regelmäßige Stürze mit hohem Verletzungsrisiko
    \item Zunehmender Unterstützungsbedarf bei der Selbstversorgung
\end{itemize}

Ärztliche Unterlagen werden nachgereicht.

Bitte teilen Sie mir den Termin für die MDK-Begutachtung rechtzeitig mit.

Mit freundlichen Grüßen\\
{[Unterschrift Bevollmächtigte/r]}
\end{graybox}

\subsection{Widerspruch bei Ablehnung}

\begin{graybox}
\small
An: [Pflegekasse]\\
Datum: [innerhalb 1 Monat nach Bescheid]\\
Betreff: \textbf{WIDERSPRUCH} gegen Bescheid vom [Datum]\\
Versicherte: [Name], geb. [Datum], Vers.-Nr.: [Nummer]\\
Aktenzeichen: [aus Bescheid]

\vspace{0.3cm}
Sehr geehrte Damen und Herren,

gegen den oben genannten Bescheid lege ich hiermit fristgerecht

\begin{center}
\textbf{\Large WIDERSPRUCH}
\end{center}

ein.

\textbf{Begründung:}\\
Die Einstufung in Pflegegrad [X] entspricht nicht dem tatsächlichen Pflegebedarf. [Begründung folgt / wird nachgereicht]

Ich bitte um Übersendung des MDK-Gutachtens sowie um Überprüfung der Einstufung.

Mit freundlichen Grüßen\\
{[Unterschrift]}
\end{graybox}

%%%%%%%%%%%%%%%%%%%%%%%%%%%%%%%%%%%%%%%%%%%%%%%%%%%%%%%%%%%%%%%%%%%%%%%%%%%%%%%
\section{Checkliste zum Ausdrucken}

\begin{graybox}
\begin{center}
\textbf{\Large CHECKLISTE: HÖHERSTUFUNG PFLEGEGRAD}
\end{center}

\vspace{0.5cm}
\textbf{PHASE 1: SOFORT}
\begin{itemize}[noitemsep]
    \item[$\square$] Höherstufungsantrag abgeschickt \hfill Datum: \_\_\_\_\_\_\_\_\_\_\_
    \item[$\square$] Sturzprotokoll begonnen \hfill Datum: \_\_\_\_\_\_\_\_\_\_\_
    \item[$\square$] Pflegeheim 1: \_\_\_\_\_\_\_\_\_\_\_\_\_\_\_\_ \hfill Warteliste: $\square$
    \item[$\square$] Pflegeheim 2: \_\_\_\_\_\_\_\_\_\_\_\_\_\_\_\_ \hfill Warteliste: $\square$
    \item[$\square$] Pflegeheim 3: \_\_\_\_\_\_\_\_\_\_\_\_\_\_\_\_ \hfill Warteliste: $\square$
    \item[$\square$] Hausnotruf installiert \hfill Datum: \_\_\_\_\_\_\_\_\_\_\_
\end{itemize}

\vspace{0.5cm}
\textbf{PHASE 2: VORBEREITUNG}
\begin{itemize}[noitemsep]
    \item[$\square$] Pflegetagebuch geführt (mind. 2 Wochen)
    \item[$\square$] Hausarzt-Attest \hfill $\square$ erhalten
    \item[$\square$] Neurologe/Psychiater-Attest \hfill $\square$ erhalten
    \item[$\square$] Medikamentenliste aktuell \hfill $\square$ ja
    \item[$\square$] Hilfsmittel dokumentiert \hfill $\square$ ja
    \item[$\square$] Schwester für MDK-Termin: \_\_\_\_\_\_\_\_\_\_\_\_\_\_\_\_
    \item[$\square$] Ehemann vorbereitet \hfill $\square$ ja
\end{itemize}

\vspace{0.5cm}
\textbf{PHASE 3: MDK-TERMIN}
\begin{itemize}[noitemsep]
    \item[$\square$] Termin \hfill Datum: \_\_\_\_\_\_\_\_\_\_\_
    \item[$\square$] Pflegetagebuch dabei \hfill $\square$ ja
    \item[$\square$] Alle Unterlagen dabei \hfill $\square$ ja
    \item[$\square$] Eigenes Protokoll geschrieben \hfill $\square$ ja
\end{itemize}

\vspace{0.5cm}
\textbf{PHASE 4: NACH BESCHEID}
\begin{itemize}[noitemsep]
    \item[$\square$] Bescheid erhalten \hfill Datum: \_\_\_\_\_\_\_\_\_\_\_
    \item[$\square$] Ergebnis: PG \_\_\_
    \item[$\square$] Bei Ablehnung: Widerspruch \hfill Datum: \_\_\_\_\_\_\_\_\_\_\_
\end{itemize}
\end{graybox}

%%%%%%%%%%%%%%%%%%%%%%%%%%%%%%%%%%%%%%%%%%%%%%%%%%%%%%%%%%%%%%%%%%%%%%%%%%%%%%%
\section{Wichtige Kontakte \& Ressourcen}

\subsection{Kostenlose Beratung in Hessen}

\begin{center}
\begin{tabular}{lll}
\toprule
\textbf{Stelle} & \textbf{Kontakt} & \textbf{Leistung} \\
\midrule
Pflegestützpunkte Hessen & pflegestuetzpunkte-hessen.de & Neutrale Beratung \\
Alzheimer Gesellschaft Hessen & 069-67 77 66 & Demenz-Beratung \\
VdK Hessen & 069-714002-0 & Hilfe bei Widerspruch \\
Compass Pflegeberatung & 0800-101 88 00 & Private Pflegevers. \\
\bottomrule
\end{tabular}
\end{center}

%%%%%%%%%%%%%%%%%%%%%%%%%%%%%%%%%%%%%%%%%%%%%%%%%%%%%%%%%%%%%%%%%%%%%%%%%%%%%%%
\section{Quellen}

\subsection{Offizielle Quellen}
\begin{itemize}[noitemsep]
    \item Bundesgesundheitsministerium -- Pflegegrade: \url{https://www.bundesgesundheitsministerium.de}
    \item Bundesgesundheitsministerium -- Verhinderungspflege: \url{https://www.bundesgesundheitsministerium.de/verhinderungspflege.html}
\end{itemize}

\subsection{Pflegekassen \& Verbände}
\begin{itemize}[noitemsep]
    \item BARMER -- Kombinationsleistung: \url{https://www.barmer.de/unsere-leistungen/pflege/kombinationsleistung-1003714}
    \item AOK -- Verhinderungspflege: \url{https://www.aok.de/pk/alltagshilfe-pflegende-angehoerige/verhinderungspflege/}
    \item Techniker Krankenkasse -- Kombinationsleistung: \url{https://www.tk.de}
    \item vdek Hessen -- Eigenanteile 2026: \url{https://www.vdek.com/LVen/HES/}
\end{itemize}

\subsection{Regionale Quellen Hessen}
\begin{itemize}[noitemsep]
    \item Hessenschau -- Pflegeheim-Kosten 2026: \url{https://www.hessenschau.de}
\end{itemize}

\subsection{Pflegeportale}
\begin{itemize}[noitemsep]
    \item Pflege.de -- Verhinderungspflege, Kombinationsleistung, 24h-Pflege: \url{https://www.pflege.de}
    \item Betanet -- Ersatzpflege: \url{https://www.betanet.de}
\end{itemize}

\subsection{24h-Pflege}
\begin{itemize}[noitemsep]
    \item Lumira Pflege -- Kosten 2026: \url{https://lumira-pflege.de}
    \item Mecasa -- Polnische Pflegekraft Kosten: \url{https://www.mecasa.de}
\end{itemize}

%%%%%%%%%%%%%%%%%%%%%%%%%%%%%%%%%%%%%%%%%%%%%%%%%%%%%%%%%%%%%%%%%%%%%%%%%%%%%%%
\vfill
\begin{center}
\rule{0.8\textwidth}{0.4pt}

\vspace{0.5cm}
\textbf{ENDE DES BERICHTS}

\vspace{0.3cm}
Bei Fragen: Diesen Bericht als Grundlage für Beratungsgespräche\\
bei Pflegestützpunkt oder Pflegekasse nutzen.

\vspace{0.5cm}
Erstellt: 25. Januar 2026\\
Session: EBF-S-2026-01-25-HLT-001

\vspace{0.3cm}
{\small Erstellt mit dem Evidence-Based Framework (EBF)}
\end{center}

\end{document}
